\chapter{Hemolytic Uremic Syndrome (HUS)}
\index{Hemolytic Uremic Syndrome (HUS)}
\label{sec:Hemolytic Uremic Syndrome (HUS)}

\section{Overview}
Hemolytic uremic syndrome (HUS) is a thrombotic microangiopathy characterized by the triad of microangiopathic hemolytic anemia, thrombocytopenia, and acute kidney injury. It is the most common cause of acute renal failure in children, most frequently following infection with Shiga toxin-producing Escherichia coli (STEC), particularly serotype O157:H7. Atypical HUS (aHUS) involves uncontrolled complement activation.
\section{Classification}

\begin{description}[font=\normalfont\bfseries,leftmargin=3em,labelwidth=2.5em]
  \item[Cause] Infectious / Genetic
  \item[Organ System] Kidneys
  \item[Pathophysiology] Endothelial injury leading to microvascular thrombosis, microangiopathic hemolytic anemia, and renal injury
  \item[Duration] Acute (days to weeks) or Chronic (atypical)
  \item[Onset] Sudden
  \item[Mode of Transmission] Foodborne (STEC-HUS); Genetic (aHUS)
  \item[Clinical Course] Acute, self-limited or chronic, relapsing
  \item[Severity] Moderate to Severe
  \item[Extent of Disease] Systemic
  \item[Age Group] Children (typical); All ages (atypical)
  \item[Epidemiology] Incidence 2--3 per 100,000 in children <5 years
  \item[ICD-11 Code] 3A21.2
\end{description}

\section{Causes}
Typical (STEC) HUS is caused by Shiga toxin-producing Escherichia coli infection, usually acquired through contaminated food or water. Shiga toxin enters the circulation and damages endothelial cells, triggering microvascular thrombosis in the kidneys and other organs. Atypical HUS involves genetic mutations in complement regulatory proteins (CFH, CFI, MCP, CFB, C3) leading to uncontrolled alternative complement pathway activation.

\section{Risk Factors}

\begin{itemize}[noitemsep]
  \item Age under 5 years (typical HUS)
  \item Consumption of undercooked ground beef or contaminated produce
  \item Exposure to contaminated water or unpasteurized dairy
  \item Contact with farm animals
  \item Genetic mutations in complement genes (aHUS)
  \item Family history of aHUS
  \item Pregnancy or postpartum period (aHUS trigger)
\end{itemize}

\section{Signs and Symptoms}

\subsection{Early symptoms}
\begin{itemize}[noitemsep]
  \item Early manifestations of hemolytic uremic syndrome (hus) may be subtle and nonspecific
\end{itemize}

\subsection{Common symptoms}
\begin{itemize}[noitemsep]
  \item \subsection{Early symptoms}
  \item \subsection{Common symptoms}
  \item \textbf{Common symptoms:}
  \item Prodrome of bloody diarrhea, abdominal pain, and vomiting (typical HUS)
  \item Pallor and fatigue from hemolytic anemia
  \item Oliguria or anuria from acute kidney injury
  \item Bruising and petechiae from thrombocytopenia
  \item Neurologic symptoms (seizures, altered mental status, stroke) in severe cases
  \item Hypertension and fluid overload
  \item Pain or discomfort in the affected area
  \item Difficulty performing daily activities
\end{itemize}

\subsection{Less common symptoms}
\begin{itemize}[noitemsep]
  \item Atypical or less common presentations of hemolytic uremic syndrome (hus) may occur
\end{itemize}

\subsection{Emergency warning signs}
\begin{itemize}[noitemsep]
  \item Severe or worsening symptoms of hemolytic uremic syndrome (hus) require immediate medical evaluation
\end{itemize}
\section{Progression and Stages}
In typical STEC-HUS, symptoms develop approximately 5--10 days after onset of diarrheal illness. The acute phase lasts 1--3 weeks, during which renal function deteriorates. Most children recover renal function with supportive care, though some progress to end-stage renal disease. Atypical HUS follows a chronic relapsing course with progressive renal damage if untreated.

\section{Complications}

\begin{itemize}[noitemsep]
  \item End-stage renal disease requiring dialysis
  \item Chronic kidney disease and hypertension
  \item Neurologic sequelae (stroke, cognitive impairment, seizures)
  \item Colonic complications (perforation, stricture)
  \item Pancreatitis and diabetes mellitus
  \item Recurrence (aHUS, particularly after renal transplant)
  \item Reduced quality of life
  \item Emotional and psychological effects
\end{itemize}

\section{Diagnosis}

\begin{itemize}[noitemsep]
  \item Complete blood count showing microangiopathic hemolytic anemia (schistocytes on smear) and thrombocytopenia
  \item Renal function tests (elevated creatinine, BUN)
  \item Stool culture or PCR for STEC O157:H7
  \item Complement levels and genetic testing for aHUS
  \item ADAMTS13 activity to exclude thrombotic thrombocytopenic purpura
  \item Kidney biopsy in atypical cases
\end{itemize}

\section{Prevention}

\begin{itemize}[noitemsep]
  \item Cook ground beef thoroughly to an internal temperature of 160°F
  \item Wash fruits and vegetables thoroughly
  \item Avoid unpasteurized dairy products
  \item Practice good hand hygiene, especially after animal contact
  \item Go for regular check-ups so any problems are caught early
\end{itemize}

\section{Treatment Options}

\begin{itemize}[noitemsep]
  \item Supportive care with fluid and electrolyte management
  \item Renal replacement therapy (dialysis) for severe acute kidney injury
  \item Plasma exchange (plasmapheresis) for severe or atypical HUS
  \item Eculizumab (complement inhibitor) for atypical HUS
  \item Blood transfusion for severe anemia
  \item Avoid antibiotics in STEC-HUS (may increase Shiga toxin release)
  \item Lifestyle changes and self-care
\end{itemize}

\section{Living with It}

\begin{itemize}[noitemsep]
  \item Follow your treatment plan as prescribed
  \item Keep up with regular appointments with your healthcare provider
  \item Adjust your daily habits as needed
  \item Reach out to family, friends, or support groups for help
  \item Keep learning about your condition and how to manage it
\end{itemize}

\section{Outlook}
In typical STEC-HUS, the prognosis is good with supportive care; over 70\% of children recover renal function fully, though some have long-term sequelae. Mortality is 3--5\% in the acute phase. Atypical HUS carries a worse prognosis without treatment; with eculizumab therapy, renal outcomes have improved substantially. Long-term follow-up is essential to monitor for chronic kidney disease and hypertension.
